\documentclass[10pt]{exam}
\printanswers
\usepackage{epsfig,amssymb}
\usepackage{xcolor}
\definecolor{darkred}{rgb}{0.5,0,0}
\definecolor{darkgreen}{rgb}{0,0.5,0}
\usepackage{hyperref}
\usepackage{fullpage}
\usepackage{tikz}
\pagestyle{empty} 
\usepackage{listings}
\setlength{\parindent}{0pt} %
\setlength{\parskip}{.25cm}
\newcommand{\comment}[1]{}
\boxedpoints
\addpoints
\usepackage{amsmath}
\usepackage{algorithm2e}
\usepackage{url}
\usepackage{enumitem}
\usepackage{graphics}

\pagestyle{headandfoot}
\firstpageheader{\Large CSCE 235}{{\Large Assignment 6}}{\Large Spring 2026}

\begin{document}

\hrule

\vspace{.50cm} Name:
\text{Charlie Pyle} 
\\ NUID: 22821006
\vspace{.50cm}

\textbf{Instructions} Follow instructions \emph{carefully}; failure to do so will result in points being deducted. 

\begin{questions}


\question[8] Prove, using induction, that  
\[
1^4 + 2^4 + 3^4 + \cdots + n^4 = \frac{n(n+1)(2n+1)(3n^2+3n-1)}{30}
\]  
whenever \(n\) is a positive integer.
\begin{parts}
 \part State and prove the basis step.
 \part State the inductive hypothesis.
 \part State the inductive conclusion.
 \part Prove the inductive conclusion by the method of induction. You must justify the relevant steps.
\end{parts}

\begin{solution}
    Enter your solution here:\\
    \textbf{A:} Base case $n = 1$:
    \[1^4 = \frac{1(1+1)(2\cdot1+1)(3\cdot1^2+3\cdot1-1)}{30}\]
    \[\sum_{n=1}^{1} n^4 = 1 = 1 \checkmark \qquad \frac{1 \cdot 2 \cdot 3 \cdot 5}{30} = \frac{30}{30} = 1 \checkmark\]
    
    \textbf{B:} Inductive hypothesis
    \[1^4 + 2^4 + 3^4 + \cdots + k^4 = \frac{k(k+1)(2k+1)(3k^2+3k-1)}{30}\]
    
    \textbf{C:} Inductive conclusion:
    \[\underbrace{1^4 + 2^4 + 3^4 + \cdots + k^4}_{\text{IH}} + (k+1)^4
    = \frac{(k+1)(k+2)(2k+3)(3(k+1)^2+3(k+1)-1)}{30}\]
    
    \textbf{D:} Prove:
    \[\frac{k(k+1)(2k+1)(3k^2+3k-1)}{30} + (k+1)^4
    = \frac{(k+1)(k+2)(2k+3)(3(k+1)^2+3(k+1)-1)}{30}\]
    
    \[\frac{k(k+1)(2k+1)(3k^2+3k-1) + 30(k+1)^4}{30}\]
    
    \[\frac{k+1}{30}\left[k(2k+1)(3k^2+3k-1) + 30(k+1)^3\right]\]
    
    Expanding $(2k^2+k)(3k^2+3k-1)$:
    \[6k^4 + 6k^3 - 2k^2 + 3k^3 + 3k^2 - k = 6k^4 + 9k^3 + k^2 - k\]
    
    And $30(k+1)^3 = 30k^3 + 90k^2 + 90k + 30$
    
    Combining:
    \[\frac{k+1}{30}\left[6k^4 + 39k^3 + 91k^2 + 89k + 30\right]\]
    
    We now expand the right-hand side target $(k+2)(2k+3)(3k^2+9k+5)$ to verify it equals the polynomial above.
    
    First, $(k+2)(2k+3) = 2k^2 + 7k + 6$.
    
    Then $(2k^2 + 7k + 6)(3k^2 + 9k + 5)$:
    \[= 6k^4 + 18k^3 + 10k^2 + 21k^3 + 63k^2 + 35k + 18k^2 + 54k + 30\]
    \[= 6k^4 + 39k^3 + 91k^2 + 89k + 30 \checkmark\]
    
    Since $3(k+1)^2 + 3(k+1) - 1 = 3k^2 + 9k + 5$, we conclude:
    \[\frac{k+1}{30}(k+2)(2k+3)(3k^2+9k+5)
    = \frac{(k+1)(k+2)(2k+3)\bigl(3(k+1)^2+3(k+1)-1\bigr)}{30}\]
\end{solution}

\question[8] Prove the following using induction whenever \( n \) is a positive integer.  
\[
\frac{1}{1 \cdot 2} + \frac{1}{2 \cdot 3} + \cdots + \frac{1}{n(n+1)} = 1 - \frac{1}{n+1}.
\]
\begin{parts}
 \part State and prove the basis step.
 \part State the inductive hypothesis.
 \part State the inductive conclusion.
 \part Prove the inductive conclusion by the method of induction. You must justify the relevant steps.
\end{parts}

\begin{solution}
    Enter your solution here:\\
    \textbf{A:} Base case $n = 1$:
    \[\frac{1}{1(1+1)} = 1 - \frac{1}{1+1} \qquad \frac{1}{2} = \frac{1}{2} \checkmark\]
    
    \textbf{B:} Inductive hypothesis:
    \[\frac{1}{1\cdot2} + \frac{1}{2\cdot3} + \cdots + \frac{1}{k(k+1)} = 1 - \frac{1}{k+1}\]
    
    \textbf{C:} Inductive conclusion:
    \[\underbrace{\frac{1}{1\cdot2} + \frac{1}{2\cdot3} + \cdots + \frac{1}{k(k+1)}}_{\text{IH}}+ \frac{1}{(k+1)(k+2)} = 1 - \frac{1}{k+2}\]
    
    \textbf{D:} Prove:
    \[1 - \frac{1}{k+1} + \frac{1}{(k+1)(k+2)}\]
    
    \[= 1 + \frac{1}{k+1}\left(-1 + \frac{1}{k+2}\right)\]
    
    \[= 1 + \frac{1}{k+1}\left(\frac{-(k+2)+1}{k+2}\right)\]
    
    \[= 1 + \frac{1}{k+1}\left(\frac{-(k+1)}{k+2}\right)\]
    
    \[= 1 - \frac{1}{k+2}\]

\end{solution}

\question[8] Prove, using induction, that 2 divides $n^3 + 3n$ whenever $n$ is a positive integer.
\begin{parts}
 \part State and prove the basis step.
 \part State the inductive hypothesis.
 \part State the inductive conclusion.
 \part Prove the inductive conclusion by the method of induction. You must justify the relevant steps.
\end{parts}

\begin{solution}
    Enter your solution here:\\
    \textbf{A:} Show $2 \mid n^3 + 3n$.
    
    Base case $n = 1$: $2 \mid 1^3 + 3(1) = 4 = 2 \cdot 2$ \checkmark
    
    \textbf{B:} Inductive hypothesis: Assume $2 \mid k^3 + 3k$.
    
    \textbf{C:} Show: $2 \mid (k+1)^3 + 3(k+1)$.
    
    \textbf{D:} Expanding:
    \[(k+1)^3 + 3(k+1) = k^3 + 3k^2 + 3k + 1 + 3k + 3\]
    
    \[= (3k^2 + 6k + 4) + \underbrace{k^3 + 3k}_{2M}\]
    
    \[= 2\!\underbrace{\left(\tfrac{3}{2}k^2 + 3k + 2\right)}_{a} + 2M\]
    
    \[= 2a + 2M = 2(a + M) = 2C\]
\end{solution}

\question[8]  An employee joined a company in 2020 with a starting salary of \$95,000. Every year, this employee receives a raise of \$1000 plus 3\% of the salary of the previous year.

\begin{parts}
\part Set up a recurrence relation for this employee's salary $n$ years after 2020. \textbf{Justify} your answer.
\part  Using the iterative method, solve the recurrence relation for this employee's salary $n$ years after 2020.
\part  Determine the salary of this employee in 2026.
\end{parts}

\begin{solution}
    Enter your solution here:\\
    \begin{align*}
    P_0 &= 95000 \\
    P_1 &= 95000(1.03) + 1000 \\
    P_2 &= [95000(1.03) + 1000](1.03) + 1000 \\
    P_3 &= P_2(1.03) + 1000
    \end{align*}
    
    \[P_n = P_{n-1}(1.03) + 1000\]
    
    \[P_n = 1.03^n(95000) + 1000\left[\frac{1.03^n - 1}{0.03}\right]\]
    
    \[P_6 = 1.03^6(95000) + 1000\left[\frac{1.03^6 - 1}{0.03}\right]\approx 119{,}903\]
\end{solution}



\question[12] For the following recurrences, provide an asymptotic characterization using the \textbf{Master theorem}. 

Note: Clearly show the values of the constants of the recurrence relation when using the Master theorem, and determine which case of the Master theorem applies to this recurrence relation.

\begin{parts}
  \part $T(n) = 2T(n/3) + 4$
  \part $T(n) = 3T(n/3) + 3n$
  \part $T(n) = 2T(n/2) + 3n^2$
  \part $T(n) = 8T(n/2) + n^3$
\end{parts}

\begin{solution}
    Enter your solution here:\\
    \textbf{A:} $a = 2,\quad b = 3,\quad d = 0$
    \[2 > 3^0 = 1\]
    \[T(n) \in \mathcal{O}\!\left(n^{\log_3 2}\right)\]
    
    \textbf{B:} $a = 3,\quad b = 3,\quad d = 1$
    \[3 = 3^1\]
    \[T(n) \in \mathcal{O}(n \log n)\]
    
    \textbf{C:} $a = 2,\quad b = 2,\quad d = 2$
    \[2 < 2^2 = 4\]
    \[T(n) \in \mathcal{O}(n^2)\]
    
    \textbf{D:} $a = 8,\quad b = 2,\quad d = 3$
    \[8 = 2^3\]
    \[T(n) \in \theta(n^3 \log n)\]
\end{solution}

\end{questions}

\end{document}